%% kapa: exam=1 round=21 year=2010 period=2 form=A subject=appraisal_law subject_ko=부동산관계법규 seq=23 number=63
%% source: Q-Net 감정평가사 기출문제 (cst00309) artlSeq=5024181 file=extracted/1008/5024181/2010년 제21회 감정평가사 1차시험 기출문제/2교시.hwp
%% conversion: kordoc 4.24.51 -> tools/gen1.py (auto); review=proofread-gpt-sol; flags=none
\begin{question}{63}
부동산 가격공시 및 감정평가에 관한 법령상 감정평가서에 관한 설명으로 옳지 않은 것은? (단, 다툼이 있으면 판례에 의함)
\begin{choices}
  \item 감정평가업자는 감정평가서의 원본과 그 관련서류를 국토해양부령이 정하는 기간 이상 보존하여야 한다.
  \item 감정평가서에는 감정평가업자의 사무소 또는 법인의 명칭을 기재하고, 감정평가를 행한 감정평가사가 그 자격을 표시하고 서명ㆍ날인하여야 한다.
  \item 표준지공시지가의 결정과 관련한 감정평가서에는 평가원인을 구체적으로 특정할 필요는 없지만 각 요인별 참작 내용과 정도가 객관적으로 납득이 갈 수 있을 정도로 설명되어야 한다.
  \item 감정평가 의뢰인이 감정평가서를 분실하여 재교부를 신청한 경우, 감정평가업자는 정당한 사유가 있는 경우를 제외하고는 감정평가서를 재교부하여야 한다.
  \item 감정평가서에 거래선례나 평가선례, 거래사례비교법, 원가법 및 수익환원법 등을 모두 공란으로 둔 채, 그 토지의 전년도 공시지가와 세평가격 및 인근 표준지의 감정가격만을 참고가격으로 삼아 별다른 설명 없이 평가의견을 추상적으로만 기재하는 것은 토지의 적정가격을 반영한 것으로 인정하기 어렵다.
\end{choices}
\end{question}
\answer{3}
